\documentclass[a4paper]{article}
\usepackage[utf8]{inputenc}
\usepackage[T1]{fontenc}
\usepackage[finnish]{babel}
\usepackage[pdftex,colorlinks]{hyperref}
\usepackage{lmodern}
\usepackage{mathabx}
\newtheorem{lause}{Lause}
\newtheorem{tehtava}{Tehtävä}
\begin{document}
\title{Tietojenkäsittelyopin teoriaa}
\author{Mikko Nummelin}
\maketitle
\tableofcontents

% Johdanto
\section{Johdanto}
Tämän kurssin tarkoitus on tarjota yleissivistävä perusta tietojenkäsittelyopin
teorialle käymällä läpi ns. lukuteoreettiset funktiot ja $\lambda$ (lambda)-laskenta.
Nämä tekniikat keksittiin aikanaan 1920-1940-luvuilla sen aikaisten huomattavien
diskreetin matematiikan uranuurtajien, kuten Alan Turingin, Kurt Gödelin,
Emil Postin ja Alonso Churchin toimesta. Tuolloin huomattiin myös, että useat
eri laskennalliset tekniikat ja algoritmit kattavat käytännössä kaiken, mitä
yleensä on tietokoneen laskettavissa, niissä on pohjimmillaan samantapainen
ohjelmoitava rakenne ja että on olemassa paljon tietoteknisiä
ongelmia, jotka ovat \emph{laskennallisesti ratkeamattomia}. Lukuteoreettiset
funktiot ja $\lambda$-laskenta ovat kaksi esimerkkiä Turing-täydellisistä
algoritmiformuloinneista, mutta näitä on lukuisia muitakin, kuten
Turing-Post-koneet, kieliopit ja muut merkkijonojen
uudelleenkirjoitusjärjestelmät.

Kurssin materiaalin läpikäymisen jälkeen on tarkoitus ymmärtää paremmin
monien nykyaikaisten ohjelmointikielten taustaa ja rakenteellista pohjaa.

% Lukuteoreettiset funktiot
\section{Lukuteoreettiset funktiot}
Lukuteoreettiset funktiot käsittelevät funktioita luonnollisilta luvuilta
luonnollisille luvuille. Ns. Gödel-numeroinnin avulla on mahdollista kuvata
myös muita funktioita numeroituvilta joukoilta numeroituville joukoille, kuten
merkkijonojen funktioita tai lukulistoja. Tämä mahdollistaa myös matemaattisen
logiikan kaavojen analysoinnin, jossa kaavat voidaan itsessään kuvata luvuiksi
ja niillä voidaan muodostaa kaavoja, jotka tutkivat itseään tai toisia kaavoja.
Tällöin päädytään myös ns. lukuteorian epäsäännöllisyyslauseeseen, jonka
mukaan on mahdotonta kehittää algoritmia, joka tutkisi mielivaltaisesta
lukuteoreettisesta funktiosta, onko sillä nollakohtia vai ei.

% P.r. funktiot
\subsection{Primitiivirekursiiviset funktiot}
Tässä osiossa käsitellään ns. primitiivirekursiiviset funktiot, jotka muodostavat
\emph{melkein} kaikki laskettavissa olevat funktiot. P.r. funktioiden
ulkopuolelle jäävät laskettavissa olevat eli lukuteoreettiset funktiot ovat
sellaisia, että niissä on potentiaalisia ikisilmukoita, eli jos funktion
kirjoittaisi ohjelmaksi, ei voida suoraan asettaa ylärajaa yksittäisten
silmukoiden suoritusmäärälle. Nämä käsitellään omassa osiossaan.

\subsubsection{Initiaalifunktiot}
Vakiofunktio $\zeta()=0$, jolla ei ole yhtään parametria on primitiivirekursiivinen.
Tämä vastaa käytännössä lukua $0$. Seuraajafunktio $\sigma(n)=n+1$ on p.r.
Tällä voidaan muodostaa kaikki luonnolliset luvut $\sigma(\zeta())=1,
\sigma(\sigma(\zeta()))=2$, jne. Lisäksi, jos meillä on lista lukuja
$(x_0,x_1,x_2,...,x_{n-1})$, voidaan tätä listaa merkitä symbolilla $\vec{x}$, ja
\emph{projisoida} siitä haluttu alkio siten, että $\pi_m^n(\vec{x})=x_m$.
Projektiot $\pi_m^n$ ovat p.r.
\begin{tehtava}
Osoita, että luku $3$ on primitiivirekursiivinen vakiofunktio ilman parametreja.
\end{tehtava}
\begin{tehtava}
Laske $\pi_1^3(5,7,12)$.
\end{tehtava}

\subsubsection{Yhdistäminen ja primitiivirekursio}
Edellisessä kappaleessa emme olleet matemaattisen tarkkoja siitä, mitä
tarkoittaa funktioiden yhdistäminen, käyttäessämme sellaisia ilmaisuja, kuten
$\sigma(\sigma(\zeta())$. Muodollisesti funktio $f$ on saatu \emph{yhdistämällä}
funktioista $g$ ja $\vec{h}$, mikäli $f(\vec{x})=g(\vec{h}(\vec{x}))$. Tällöin
voidaan muodollisesti merkitä $f=(g\vec{h})$.
Lisäksi on huomattava, että funktio $f$ voidaan määritellä funktioiden $g$ ja $h$
avulla primitiivirekursiolla:

\begin{equation}
\left\{
\begin{array}{lcl}
f(0,\vec{x}) & = & g(\vec{x}) \\
f(n+1,\vec{x}) & = & h(n, f(n,\vec{x}), \vec{x})
\end{array}
\right.
\end{equation}

Tällöin voidaan muodollisesti merkitä, että $f=(\rho gh)$. Näitä muodollisia
merkintöjä tarvitaan lähinnä siinä vaiheessa, kun todistetaan lauseita, joissa
p.r. tai lukuteoreettisilla funktioilla viitataan toisiin p.r. tai lukuteoreettisiin
funktioihin.

\subsubsection{Eräitä primitiivirekursiivisia funktioita}
Olkoon $C_m^n$ $n$-argumenttinen vakiofunktio, jonka arvo on $m$. Vakiofunktiot
ovat primitiivirekursiivisia:

\begin{equation}
\left\{
\begin{array}{lcl}
C_0^0 & = & \zeta()=0 \\
C_{m+1}^0 & = & \sigma C_{m}^0 \\
C_m^{n+1} & = & C_m^n
\end{array}
\right.
\end{equation}

Summa $x+y$ on primitiivirekursiivinen:

\begin{equation}
\left\{
\begin{array}{lcl}
x + 0 & = & x \\
x + (y+1) & = & \sigma(x + y)
\end{array}
\right.
\end{equation}

\emph{Rajoitettu edeltäjä} $\sigma^{-1}$ on primitiivirekursiivinen:

\begin{equation}
\left\{
\begin{array}{lcl}
\sigma^{-1}(0) & = & 0 \\
\sigma^{-1}(n+1) & = & n
\end{array}
\right.
\end{equation}

Tätä sanotaan rajoitetuksi edeltäjäksi siksi, että luvulla $0$ ei varsinaisesti
ole edeltäjää luonnollisten lukujen joukossa ja siksi senkin rajoitettu edeltäjä
on määritelty $0$:ksi.

\emph{Rajoitettu vähennyslasku} $x \dotdiv y$ on primitiivirekursiivinen:

\begin{equation}
\left\{
\begin{array}{lcl}
x \dotdiv 0 & = & x \\
x \dotdiv (y+1) & = & \sigma^{-1}(x \dotdiv y)
\end{array}
\right.
\end{equation}

Tätä sanotaan rajoitetuksi vähennyslaskuksi siksi, että $x \dotdiv y=0$, jos $y \ge x$.
Tämä johtuu siitä, että käsittelemme vain luonnollisia lukuja emmekä negatiivisia
lukuja.

% Mu-rekursiiviset funktiot
\subsection{$\mu$-rekursiiviset funktiot}

% Lambda-laskenta
\section{$\lambda$-laskenta}
\end{document}
